% Generated by scripts/tikz_reviews.py using the compiled book's labels.
\pdfbookmark[0]{Chapter 17}{comparison-chapter-17}
\ReviewFigure{fig-deep-fc-cnn}
{A fully connected feedforward network}
{figures/deep-learning/fc.png}
{figures/tikz/deep-learning-fully-connected.pdf}
{Deep MLP layers are affine maps followed by nonlinear activations.}
{Retains the original three inputs, two hidden layers of four units, and three outputs. Every adjacent-layer pair is connected; line crossings are not nodes. The layer vectors are x0, x1, x2, x3, and component labels follow that notation. The output x3 denotes scores, avoiding confusion with the chapter target label y; a separate logistic map converts scores into probabilities.}
{deep-learning-fully-connected}
{Panel 1 of 2}
{17.1}
\ReviewFigure{fig-deep-fc-cnn}
{Feature maps in a convolutional network}
{figures/deep-learning/cnn.png}
{figures/tikz/deep-learning-convolutional.pdf}
{Convolutional layers arrange activations by spatial location and channel; downsampling reduces spatial resolution.}
{Interpretation to check: The low-resolution dimension labels are replaced by the chapter notation: n\_l=bar n\_l*d\_l and RGB input d\_0=3. The five tensor blocks show an input, convolutional features, their downsampled version, a second feature stack, and a second downsampling. The activation arrows use the pointwise tilde-rho symbols because the chapter rho already includes downsampling. The output is the score vector x\_L. Depth illustrates channel count and face size illustrates spatial resolution; exact numerical widths or filter sizes cannot be recovered and are not invented.}
{deep-learning-convolutional}
{Panel 2 of 2}
{17.1}
